\chapter{Additional Experimental Figures}
\label{app:additional-results}

Chapter~\ref{chap:results} uses seven figures to develop the main argument.
The figures below retain the other matched comparisons from the same
simulation run. Each caption gives its population and sample settings,
and each interpretation concerns that figure alone. The supplementary
results report identified in Appendix~\ref{app:evidence-index} contains
the remaining sample sizes and numerical results for every plotted point.

\section{Remaining square-table comparisons}

The principal design crosses four square table sizes with three margin
profiles. Chapter~\ref{chap:results} discusses the uniform \(2\times2\)
and \(8\times8\) comparisons and the different-skew \(8\times8\)
comparison. The nine other shape--margin combinations follow in the
same order as the design grid.

\mainlandscapefigure{2x2}{\(2\times2\)}{same_skew}{both \(P\) and \(Q\) have row and column margins \(m_2(0.8)=(0.8,0.2)\)}{same skew}{The null points are below 0.05, especially for Expanded Welch. At \(n=10\), many Expanded Welch calculations are invalid (hollow markers), so its low unconditional rate has more than one cause. Detection rises with the constructed difference mainly at \(n=1{,}000\); there Wald remains slightly higher. This skewed pair gives no general calibration advantage to the more conservative reference.}
\mainlandscapefigure{2x2}{\(2\times2\)}{different_skew}{\(P\) has row and column margins \(m_2(0.7)=(0.7,0.3)\), while \(Q\) has \(m_2(0.8)=(0.8,0.2)\)}{different skew}{At zero difference, the two tables have equal MI but different margins. Both methods are mostly below the 0.05 target, and Expanded Welch is lower. At \(n=10\), its hollow markers also indicate substantial invalidity. The alternative curves rise noticeably at \(n=1{,}000\), with Wald just above Expanded Welch. The same conservative pattern therefore persists when the equal-MI null does not require identical populations.}
\mainlandscapefigure{3x3}{\(3\times3\)}{uniform}{both \(P\) and \(Q\) have uniform row and column margins \((1/3,1/3,1/3)\)}{uniform}{At \(n=10\), Wald's zero-difference rate is above 0.05 and Expanded Welch brings it below the line; both methods are almost always valid. At \(n=100\), however, both start well below 0.05 and the correction lowers rejection further. Detection grows strongly at \(n=1{,}000\), where Wald remains slightly higher. Whether the adjustment helps calibration changes with sample size, not just table shape.}
\mainlandscapefigure{3x3}{\(3\times3\)}{same_skew}{both \(P\) and \(Q\) have row and column margins \(m_3(0.8)=(0.8,0.1,0.1)\)}{same skew}{The \(n=10\) null rate falls from slightly above 0.05 for Wald to below it for Expanded Welch, but Expanded Welch is valid in only about two thirds of those samples. At \(n=100\) and 1,000, both null rates are below 0.05. Rejection rises with the MI difference most visibly at \(n=1{,}000\); Expanded Welch remains lower. The small-sample null reduction is not solely a reference-distribution effect.}
\mainlandscapefigure{3x3}{\(3\times3\)}{different_skew}{\(P\) has row and column margins \(m_3(0.7)=(0.7,0.15,0.15)\), while \(Q\) has \(m_3(0.8)=(0.8,0.1,0.1)\)}{different skew}{With different margins, Wald rejects too often at the \(n=10\) null; Expanded Welch rejects less, although its hollow markers show reduced validity. At \(n=100\), both are conservative. At \(n=1{,}000\), detection increases with the true difference and the methods draw closer. This comparison shows a small-sample benefit against liberal rejection but not a consistent improvement across sample sizes.}
\mainlandscapefigure{5x5}{\(5\times5\)}{uniform}{both \(P\) and \(Q\) have uniform row and column margins \((1/5,\ldots,1/5)\)}{uniform}{Both methods are valid in the displayed panels. At zero difference, Wald is liberal at \(n=10\), while Expanded Welch reduces but does not eliminate the excess. At \(n=100\) and 1,000, both null rates are below 0.05. Little detection is gained over the displayed difference range at \(n=100\), whereas rejection grows clearly at \(n=1{,}000\). The correction lowers both false positives and detection.}
\mainlandscapefigure{5x5}{\(5\times5\)}{same_skew}{both \(P\) and \(Q\) have row and column margins \(m_5(0.8)=(0.8,0.05,0.05,0.05,0.05)\)}{same skew}{At \(n=10\), Wald starts above the 0.05 line and Expanded Welch below it, but the latter is invalid for roughly one third of samples. At \(n=100\) and 1,000, both begin below the line. Detection increases substantially only in the largest-sample panel, with Wald higher. The small-sample appearance of better calibration must be read together with validity.}
\mainlandscapefigure{5x5}{\(5\times5\)}{different_skew}{\(P\) has row and column margins \(m_5(0.7)=(0.7,0.075,0.075,0.075,0.075)\), while \(Q\) has \(m_5(0.8)=(0.8,0.05,0.05,0.05,0.05)\)}{different skew}{The \(n=10\) null is liberal for both methods, though less so for Expanded Welch; its hollow markers also show reduced validity. At \(n=100\), rejection stays low and even falls across this short alternative range, so these curves should not be described as increasing power. At \(n=1{,}000\), detection rises and Wald remains slightly higher. The benefit of reducing excess null rejection is regime-specific.}
\mainlandscapefigure{8x8}{\(8\times8\)}{same_skew}{both \(P\) and \(Q\) have row and column margins \(m_8(0.8)=(0.8,0.2/7,\ldots,0.2/7)\)}{same skew}{At \(n=10\), Wald is above 0.05 at the null and Expanded Welch below it, but many Expanded Welch outputs are invalid. Both null rates are below 0.05 at \(n=100\) and 1,000. Across the alternative points, appreciable growth appears chiefly at \(n=1{,}000\), with Wald higher. The adjustment does not turn the smaller-sample curves into strong detection.}

\FloatBarrier

\section{Other controlled comparisons}

The following figures hold the stated margins, MI targets and sample
sizes fixed while changing association location, sample allocation,
table shape or population construction. The final figure starts at
exact independence, where the first-order MI variance is zero.

\begin{figure}[htbp]
\centering
\includegraphics[width=0.96\textwidth]{patterns_5x5_different_skew.pdf}
\caption[Changed-cell comparison, \(5\times5\) tables]{Changed-cell comparison for \(5\times5\) tables. Columns place
association in the first \(2\times2\) block, last (rare) \(2\times2\)
block, or a centred score outer product spread across all cells. For every
column, \(P\)'s row and column margins are
\((0.7,0.075,0.075,0.075,0.075)\), and \(Q\)'s are
\((0.8,0.05,0.05,0.05,0.05)\). Here \(I(P)=0.0001\),
\(I(Q)=I(P)+|\Delta|\), \(n_P=n_Q=100\), and
\(|\Delta|=0,0.0001,0.00025,0.0005,0.001,0.002\) nats.
Top: 0--1; bottom: 0--0.15. Each point has 20,000 replicates.}
\label{fig:patterns-focus}
\resultinterpretation{For first, rare and spread patterns, respectively, the
zero-difference Wald rates are \(\PatternFirstWaldRate\),
\(\PatternRareWaldRate\) and \(\PatternSpreadWaldRate\); Expanded Welch gives
\(\PatternFirstExpandedRate\), \(\PatternRareExpandedRate\) and
\(\PatternSpreadExpandedRate\). All displayed calculations are valid. The
alternative curves remain nearly flat over this short MI range, so neither
method detects these differences well. Changing the location of association
does not remove the conservative gap here, although this one plot does not
make all patterns equivalent.}
\end{figure}

\begin{figure}[htbp]
\centering
\includegraphics[width=0.96\textwidth]{imbalance_3x3_q_higher.pdf}
\caption[Unequal samples, \(Q\) higher]{Unequal samples, \(Q\) higher. Columns have
\((n_P,n_Q)=(50,250)\) and \((250,50)\). For both columns, \(P\) has
\(3\times3\) row and column margins \((0.7,0.15,0.15)\) and \(Q\) has
\((0.8,0.1,0.1)\); the first-cell change block is used.
\(I(P)=0.02\) and \(I(Q)=0.02+|\Delta|\), with
\(|\Delta|=0,0.001,0.002,0.005,0.01,0.02\) nats. Top: 0--1; bottom:
0--0.15. Each point has 20,000 replicates.}
\label{fig:imbalance-q}
\resultinterpretation{At zero difference, Wald is near or above 0.05 while
Expanded Welch is below it in both allocations. When \(Q\) has more observations
\((n_P,n_Q)=(50,250)\), rejection rises with \(I(Q)-I(P)\); at 0.02 nats it
reaches \(\ImbalanceQHigherWaldRate\) for Wald and
\(\ImbalanceQHigherExpandedRate\) for Expanded Welch. With \(P\) larger, the
curves rise much less. These points are valid, but the alternative gap is not
a size-adjusted power comparison because their matched null rates differ.}
\end{figure}

\begin{figure}[htbp]
\centering
\includegraphics[width=0.96\textwidth]{imbalance_3x3_p_higher.pdf}
\caption[Unequal samples, \(P\) higher]{Unequal samples, \(P\) higher. Columns have
\((n_P,n_Q)=(50,250)\) and \((250,50)\). \(P\)'s \(3\times3\) row and
column margins are \((0.7,0.15,0.15)\); \(Q\)'s are \((0.8,0.1,0.1)\).
Both use the first-cell change block. \(I(Q)=0.02\) and
\(I(P)=0.02+|\Delta|\), with
\(|\Delta|=0,0.001,0.002,0.005,0.01,0.02\) nats.
Top: 0--1; bottom: 0--0.15. Each point has 20,000 replicates.}
\label{fig:imbalance-p}
\resultinterpretation{The zero-difference rates are the same as in
Figure~\ref{fig:imbalance-q}, because each allocation uses the same pair of
populations in both figures at that point; \(P\) and \(Q\) still have different
margins. Reversing the MI direction changes what happens next: when \(P\) has
the larger sample, Wald rises well above its null rate, whereas it stays
nearly flat when \(Q\) has the larger sample. At \((n_P,n_Q)=(50,250)\) and
0.02 nats, Wald rejects
\(\ImbalancePHigherWaldRate\) and Expanded Welch
\(\ImbalancePHigherExpandedRate\); both calculations are valid. Direction and
allocation therefore matter even at the same absolute MI difference.}
\end{figure}

\begin{figure}[htbp]
\centering
\includegraphics[width=0.94\textwidth]{rectangular_different_skew.pdf}
\caption[Rectangular tables]{Rectangular tables. Columns are \(2\times3\) and \(3\times5\).
In each dimension \(k\), \(P\)'s corresponding margin is \(m_k(0.7)\)
and \(Q\)'s is \(m_k(0.8)\), for both rows and columns. Thus the first
column uses \(P\) rows \((0.7,0.3)\), columns
\((0.7,0.15,0.15)\), and \(Q\) rows \((0.8,0.2)\), columns
\((0.8,0.1,0.1)\); the second applies the same formula to 3 and 5
categories. Both use first-cell changes, \(n_P=n_Q=100\),
\(I(P)=0.02\), and \(I(Q)=0.02+|\Delta|\) for
\(|\Delta|=0,0.001,0.002,0.005,0.01,0.02\) nats.
Top: 0--1; bottom: 0--0.15. Each point has 20,000 replicates.}
\label{fig:rectangular-focus}
\resultinterpretation{Both rectangular panels start below 0.05, with
Expanded Welch lower. Rejection grows only modestly by 0.02 nats and remains
below 0.05 for both methods throughout the displayed range; validity is high.
Thus these shapes show weak detection over the chosen differences, not an
improvement in false-positive calibration from the lower Expanded Welch
curve. The conservative pattern is not confined to square tables.}
\end{figure}

\begin{figure}[htbp]
\centering
\includegraphics[width=0.94\textwidth]{construction_5x5_different_skew.pdf}
\caption[Additive versus log-linear construction]{Additive versus ordinal log-linear construction for
\(5\times5\) tables. The left column uses the first-cell additive block;
the right fits an ordinal interaction while preserving the margins. In both,
\(P\)'s row and column margins are
\((0.7,0.075,0.075,0.075,0.075)\) and \(Q\)'s are
\((0.8,0.05,0.05,0.05,0.05)\). Both have
\(I(P)=0.02\), \(I(Q)=0.02+|\Delta|\), \(n_P=n_Q=100\), and
\(|\Delta|=0,0.001,0.002,0.005,0.01,0.02\) nats.
Top: 0--1; bottom: 0--0.15. Each point has 20,000 replicates.}
\label{fig:construction-focus}
\resultinterpretation{At zero difference, the additive and log-linear Wald
rates are \(\ConstructionAdditiveWaldRate\) and
\(\ConstructionLoglinearWaldRate\); Expanded Welch gives
\(\ConstructionAdditiveExpandedRate\) and
\(\ConstructionLoglinearExpandedRate\). All four are valid and below 0.05.
Rejection stays low or falls slightly as the displayed difference grows.
Matching margins and MI therefore does not make finite-sample behaviour
identical across constructions, but neither construction removes the
conservative gap between the methods.}
\end{figure}

\begin{figure}[htbp]
\centering
\includegraphics[width=0.96\textwidth]{independence_2x2.pdf}
\caption[Exact-independence boundary]{Exact-independence boundary for \(2\times2\) tables. Columns are
uniform margins \((0.5,0.5)\) for both populations, same skew
\((0.8,0.2)\) for both, and different skew with \(P=(0.7,0.3)\) and
\(Q=(0.8,0.2)\); each vector is used for both rows and columns. The
first-cell block is used, \(I(P)=0\), \(I(Q)=|\Delta|\),
\(n_P=n_Q=100\), and
\(|\Delta|=0,0.001,0.002,0.005,0.01,0.02\) nats.
Top: 0--1; bottom: 0--0.15. Each point has 20,000 replicates.}
\label{fig:independence-focus}
\resultinterpretation{At the uniform zero-difference point, Wald rejects
0.00025 and Expanded Welch 0.00010 of samples; validity is high. Rejection grows little over
most of the displayed alternative range, and Expanded Welch remains lower.
This is not a routine positive-MI validation: at exact independence the
population first-order MI variance is zero, so the approximation underlying
both curves changes. Appendix~\ref{app:independence-boundary} explains the
different reference used by a one-table independence test.}
\end{figure}

\FloatBarrier

\section{Main-grid null-rate overview}

Table~\ref{tab:null-band-summary} describes where the 108 main-grid
null rates fall relative to 0.025--0.075. This interval is only a
descriptive aid, not a formal calibration test. The rates and
validity fractions for each configuration are available separately
in the supplementary results report and saved data.

\begin{table}[htbp]
\centering
\small
\caption{Main-grid null configurations below, inside or above the descriptive rejection-rate interval 0.025--0.075.}
\label{tab:null-band-summary}
\begin{tabular}{@{}llrrrr@{}}
\toprule
Margins & Method & Below & Inside & Above & Validity below 90\%\\
\midrule
Uniform & Normal Wald & 10 & 18 & 8 & 4 \\
Uniform & Expanded Welch & 19 & 16 & 1 & 5 \\
Same skew & Normal Wald & 14 & 15 & 7 & 8 \\
Same skew & Expanded Welch & 25 & 11 & 0 & 12 \\
Different skew & Normal Wald & 9 & 16 & 11 & 8 \\
Different skew & Expanded Welch & 19 & 13 & 4 & 12 \\
\bottomrule
\end{tabular}
\end{table}


\FloatBarrier
